\documentclass[a4paper,11pt]{article}
\usepackage[a4paper,left=25mm,right=25mm,top=25mm,bottom=25mm]{geometry}
\usepackage[T1]{fontenc}
\usepackage[utf8]{inputenc}
\usepackage{lmodern}
\usepackage{setspace}
\usepackage{hyperref}
\usepackage{enumitem}
\usepackage{longtable}
\usepackage{booktabs}
\usepackage{array}

\setlength{\parskip}{0.45em}
\setlength{\parindent}{0pt}
\setstretch{1.15}
\renewcommand{\arraystretch}{1.15}

\title{Project Working Explanation\\\large Local PDF QA System (Theory Version)}
\author{Final Year Project Documentation}
\date{\today}

\begin{document}
\maketitle
\tableofcontents
\newpage

\section{Introduction}
This project is a local question-answering system that learns from PDF documents.
The full workflow starts from raw PDF files and ends with a command-line interface where a user asks a question and receives an answer.
The system is built on a modern, high-performance GPT-style decoder architecture capable of running efficiently on GPUs or falling back to CPUs, while remaining fully offline.
The core idea is practical:
\begin{itemize}[leftmargin=1.5em]
    \item First, try to generate an answer using the trained language model.
    \item If generated text is weak or unreliable, use retrieval fallback from the prepared sentence bank.
\end{itemize}

This hybrid behavior improves stability without requiring external APIs or cloud models.

\section{Overall Working Theory}
The project can be understood as two connected phases:
\begin{enumerate}[leftmargin=1.5em]
    \item \textbf{Knowledge preparation and training phase}
    \item \textbf{Question answering phase}
\end{enumerate}

In the first phase, the project converts unstructured PDF text into a clean training resource and trains a Transformer model.
In the second phase, the trained model is used to answer user questions with a fallback mechanism that protects against irrelevant or low-quality outputs.

\section{Phase 1: Knowledge Preparation and Training}
\subsection{Step 1: PDF Text Extraction}
The process begins in the \texttt{data/raw/} directory.
All PDF files are read and converted into text using a PDF text extraction library.
Each PDF produces one text file in \texttt{data/extracted/}.

\textbf{Theory:}
PDFs are presentation-oriented files, not structured knowledge files.
Extraction converts visual document content into machine-readable plain text so downstream NLP steps can operate on it.

\subsection{Step 2: Text Cleaning and Normalization}
Extracted text usually contains noise such as odd symbols, repeated spaces, OCR artifacts, and formatting breaks.
The cleaning module normalizes punctuation and spacing and merges all cleaned text into a single corpus file.

\textbf{Theory:}
A language model is very sensitive to input quality.
If noisy patterns are not removed, the model learns these mistakes and generates unstable answers later.
Cleaning is therefore a quality-control stage, not just formatting.

\subsection{Step 3: Sentence Segmentation}
The cleaned corpus is split into sentence-level lines and stored in \texttt{data/cleaned/sentences.txt}.

\textbf{Theory:}
Sentence-level storage has two direct benefits:
\begin{itemize}[leftmargin=1.5em]
    \item It gives retrieval fallback a compact factual memory source.
    \item It provides manageable units for creating QA-style training samples.
\end{itemize}

\subsection{Step 4: QA-style Dataset Creation}
The dataset creator converts sentence lines into multi-turn chat format using special tokens (\texttt{<|user|>}, \texttt{<|assistant|>}), generating up to 7 variants per sentence (definition, explanation, descriptive styles, etc.):
\begin{verbatim}
<|user|> What is <concept>?
<|assistant|> <sentence>
\end{verbatim}

Current implementation uses a rich data augmentation strategy:
\begin{enumerate}[leftmargin=1.5em]
    \item Primary question from sentence pattern (definition-style)
    \item Explanation and descriptive variants
    \item Multi-sentence context windows for longer answers
\end{enumerate}

\textbf{Theory:}
This step transforms passive text into supervised interaction format.
Because inference will also use the same Ask/Answer prompt style, training and inference stay aligned.

\subsection{Step 5: Tokenizer Training}
A SentencePiece tokenizer is trained on the generated training text.
It produces:
\begin{itemize}[leftmargin=1.5em]
    \item \texttt{tokenizer.model}
    \item \texttt{tokenizer.vocab}
\end{itemize}

\textbf{Theory:}
Tokenization is the bridge between language and model input.
Instead of using full words only, subword tokenization handles rare and domain-specific terms better and makes training more robust on limited data.

The training script loads tokenizer and QA records, tokenizes text, builds batches, and trains the Transformer using Mixed Precision (FP16) on GPU (or CPU). It employs an AdamW optimizer, a Cosine Learning Rate schedule with linear warmup, and gradient clipping, saving the best \texttt{model\_best.pth} based on validation loss.
\textbf{Theory:}
The model learns next-token prediction from contextual patterns.
In simple terms, it learns how answers usually continue when a question is written in the project prompt format.

\section{Model Theory}
\subsection{Architecture Choice}
The project uses a modern GPT-style decoder-only Transformer model. It incorporates advanced, state-of-the-art techniques such as RMSNorm (Pre-Norm) for training stability, Rotary Positional Embeddings (RoPE) to remove hard sequence length limits, SwiGLU Feed-Forward Networks for better expressivity, and Grouped Query Attention (GQA) for parameter and memory efficiency.

\textbf{Why Transformer:}
\begin{itemize}[leftmargin=1.5em]
    \item It captures context better than simple n-gram or shallow sequence models.
    \item It scales reasonably for small local experiments.
    \item It supports autoregressive answer generation naturally.
\end{itemize}

\subsection{Causal Behavior and KV Cache}
The causal attention mask ensures the model predicts tokens using only current and past context, not future tokens. During inference, a Key-Value (KV) Cache is utilized to store past attention states, meaning the model only needs to process the newest token rather than recomputing the entire sequence.

\textbf{Theory:}
This matches real answering behavior autoregressively and speeds up generation by 5-10x by avoiding redundant matrix multiplications.

\section{Phase 2: Question Answering Theory}
\subsection{Prompt-driven Generation}
At runtime, the user enters a question.
The system creates a prompt in the same format used during dataset creation:
\begin{verbatim}
<|user|> <user question>
<|assistant|>
\end{verbatim}
Then the model generates answer tokens using advanced sampling settings including Top-K, Nucleus (Top-P) sampling, and repetition penalties.

\textbf{Theory:}
Prompt consistency helps reduce distribution mismatch.
The model sees familiar structure, which improves coherence on small training corpora.

\subsection{Weak Answer Detection}
Generated output is checked with lightweight quality conditions such as:
\begin{itemize}[leftmargin=1.5em]
    \item too short
    \item malformed markers
    \item poor keyword coverage
\end{itemize}

\textbf{Theory:}
A local small model may produce fluent but irrelevant text.
Weak-answer detection acts as a reliability filter before final output is shown.

\subsection{Retrieval Fallback Theory}
If generation looks weak, the system searches \texttt{sentences.txt} using a TF-IDF cosine similarity approach and scores candidate lines by relevance.
Current logic uses semantic overlap plus subject-aware checks for definition-style questions.

\textbf{Key reliability rule:}
For questions like ``What is ...'', fallback retrieval is accepted only if the candidate contains the asked subject phrase.
This prevents topic drift, for example selecting a sentence from a different law or concept.

\subsection{Final Decision Layer}
The answer selection policy is intentionally simple:
\begin{itemize}[leftmargin=1.5em]
    \item Use generation when it is acceptable.
    \item Use retrieval when generation is weak and retrieval confidence is strong.
    \item If both are weak, return a safe ``not found clearly'' style response.
\end{itemize}

\textbf{Theory:}
This decision layer trades style for correctness.
It is better to return a conservative message than to return a confident but incorrect answer.

\section{Role of Each Main File}
\begin{longtable}{>{\raggedright\arraybackslash}p{0.28\linewidth} >{\raggedright\arraybackslash}p{0.67\linewidth}}
\toprule
\textbf{File} & \textbf{Role in system behavior} \\
\midrule
\texttt{textextraction.py} & Converts raw PDFs to text files. \\
\texttt{clean\_text.py} & Removes text noise and builds a clean merged corpus. \\
\texttt{sentence\_split.py} & Builds sentence-level knowledge source for retrieval and dataset input. \\
\texttt{create\_dataset.py} & Creates QA-style training records from sentence data. \\
\texttt{train\_tokenizer.py} & Learns tokenizer vocabulary from project training text. \\
\texttt{model\_def.py} & Defines the modern GPT-style Transformer network design. \\
\texttt{train.py} & Trains model with AdamW, Cosine LR, Mixed Precision, and exports checkpoint. \\
\texttt{generate.py} & Handles interactive QA with KV-cache generation and TF-IDF fallback logic. \\
\texttt{app.py} & Provides a Gradio-based interactive web interface. \\
\texttt{run\_pipeline.py} & Orchestrates full pipeline execution with simple options. \\
\bottomrule
\end{longtable}

\section{Why This Simplified Design Is Useful}
\begin{itemize}[leftmargin=1.5em]
    \item Easy to understand and explain in academic evaluation.
    \item Easy to debug because each stage is separate and script-based.
    \item Fully local and reproducible.
    \item Safer outputs due to fallback and confidence gating.
\end{itemize}

\section{Current Limitations}
\begin{itemize}[leftmargin=1.5em]
    \item Sentence-level retrieval can miss context spread across multiple lines.
    \item Small local training data limits generation fluency.
    \item Domain mismatch questions may return safe fallback response.
\end{itemize}

\section{Suggested Improvements}
\begin{enumerate}[leftmargin=1.5em]
    \item Move from sentence-level retrieval to chunk-level retrieval.
    \item Add semantic similarity scoring (e.g., dense embeddings) for better relevance.
    \item Expand the dataset with more diverse QA pairs to improve generation fluency.
\end{enumerate}

\section{How to Run}
\subsection{Train and build artifacts}
\begin{verbatim}
.\.venv\Scripts\python.exe run_pipeline.py --epochs 15
\end{verbatim}

\subsection{Ask questions via Web Interface}
\begin{verbatim}
.\.venv\Scripts\python.exe app.py
\end{verbatim}

\subsection{Ask questions via CLI}
\begin{verbatim}
.\.venv\Scripts\python.exe generate.py
\end{verbatim}

\section{Conclusion}
This project demonstrates a complete local QA workflow from raw PDFs to interactive answers.
Its strongest practical feature is the hybrid inference strategy:
model generation for flexibility and retrieval fallback for reliability.
The system is intentionally basic, clear, and suitable for final year project demonstration and further extension.

\end{document}
